\clearpage\thispagestyle{empty}
\pdfbookmark[0]{Quellen zur Besprechung der Mailänder Ephemeriden für 1814}{carlini1814-sources}
\begin{center}\Large\bfseries Quellen zur Besprechung der Mailänder Ephemeriden für 1814\end{center}

Die Übertragung folgt der Besprechung von Carl Friedrich Gauß in den \emph{Werken}, Bd.~VI (Göttingen, 1874), S.~566--567. Der Erstdruck erschien in den \emph{Göttingischen gelehrten Anzeigen}, Stück~48 vom 24.~März 1814, S.~475--478. Der Sammelband ergänzt den eingeklammerten Seitenverweis auf die Besprechung des Vorjahrs.

Das Titelblatt des Exemplars der Österreichischen Nationalbibliothek nennt Francesco Carlini als Berechner der Ephemeriden für 1814 und gibt Mailand, Stamperia Reale, 1813 an. Das Inhaltsverzeichnis schreibt die auf die Epoche 1810 reduzierten mittleren Sternpositionen Carlo Brioschi zu. Google hat das Exemplar digitalisiert. Es ist über Mario Carpinos digitale Bibliografie an der Sternwarte Brera des INAF zugänglich.

Barnaba Oriani bespricht den vermeintlichen Kometenkern im Anhang (S.~6). Giuseppe Piazzis Brief vom 4.~Juli 1812 steht auf S.~70. Die meteorologischen Tabellen von G.\,Angelo Cesaris für 1810--1812 beginnen auf S.~105, 117 und 129. Die Tabelle für Juli 1811 auf S.~123 gibt den höchsten Thermometerstand mit $+25{,}4$ an. Die Übertragung behält die Form $25^\circ\,4$ der Besprechung bei.

KI-Übertragung des deutschen Originals in heutiges Deutsch.

Erstdruck (\href{https://images.p.sub.uni-goettingen.de/iiif/image/gdz:PPN319721507_1814:00000477/full/1650,/0/default.jpg}{475}, \href{https://images.p.sub.uni-goettingen.de/iiif/image/gdz:PPN319721507_1814:00000478/full/1650,/0/default.jpg}{476}, \href{https://images.p.sub.uni-goettingen.de/iiif/image/gdz:PPN319721507_1814:00000479/full/1650,/0/default.jpg}{477}, \href{https://images.p.sub.uni-goettingen.de/iiif/image/gdz:PPN319721507_1814:00000480/full/1650,/0/default.jpg}{478}).

\href{https://www1.brera.inaf.it/~mario.carpino/Biblioteca_Digitale/Effemeridi_Milano/mirror/EM1814b.pdf#page=6}{Titelblatt des Jahrbuchs}, \href{https://www1.brera.inaf.it/~mario.carpino/Biblioteca_Digitale/Effemeridi_Milano/mirror/EM1814b.pdf#page=12}{Inhaltsverzeichnis}, \href{https://www1.brera.inaf.it/~mario.carpino/Biblioteca_Digitale/Effemeridi_Milano/mirror/EM1814b.pdf#page=139}{Oriani über den Kometenkern}, \href{https://www1.brera.inaf.it/~mario.carpino/Biblioteca_Digitale/Effemeridi_Milano/mirror/EM1814b.pdf#page=256}{Cesaris’ Julitabelle}, \href{https://images.p.sub.uni-goettingen.de/iiif/image/gdz:PPN319721507_1813:00000949/full/1650,/0/default.jpg}{Frühere Besprechung}.

\clearpage
